% ============================================================
% CONFIGURAÇÃO DO SUMÁRIO (ABNT NBR 6027 / IFSul)
% ============================================================
\cleardoublepage

% Garante que todas as páginas do Sumário não exibam numeração
\pagestyle{empty}
\thispagestyle{empty}
\addtocontents{toc}{\protect\thispagestyle{empty}}

% Título do Sumário centralizado e em caixa alta
\renewcommand{\contentsname}{%
    \vspace{-1.5cm}
    \begin{center}
        \textbf{\large SUMÁRIO}
    \end{center}
    \vspace{0.8cm}
}

% Adiciona linhas pontilhadas (leaders) também nas seções primárias
\makeatletter
\renewcommand*\l@section[2]{%
  \ifnum \c@tocdepth >\z@
    \addpenalty\@secpenalty
    \addvspace{0.8em \@plus\p@}%
    \setlength\@tempdima{1.8em}%
    \begingroup
      \parindent \z@ \rightskip \@pnumwidth
      \parfillskip -\@pnumwidth
      \leavevmode \bfseries
      \advance\leftskip\@tempdima
      \hskip -\leftskip
      #1\nobreak\normalfont\leaders\hbox{$\m@th
        \mkern \@dotsep mu\hbox{.}\mkern \@dotsep
        mu$}\hfill\nobreak\hb@xt@\@pnumwidth{\hss \bfseries #2}\par
    \endgroup
  \fi}
\makeatother

% Bookmark clicável para o Sumário nos visualizadores de PDF
\pdfbookmark[0]{Sumário}{toc}

% Renderização do Sumário
\tableofcontents

\clearpage
